\begingroup
\setlength{\LTcapwidth}{\textwidth}
\begin{longtable}{p{0.32\textwidth}p{0.53\textwidth}}
\caption{Struktur metadata dan evidence pada basis pengetahuan.}
\label{tab:tab5.8_struktur_metadata_dan_evidence} \\
\toprule
\textbf{Elemen} & \textbf{Fungsi} \\
\midrule
\endfirsthead

\multicolumn{2}{@{}l}{\small\emph{Tabel \thetable{} (lanjutan)}} \\
\toprule
\textbf{Elemen} & \textbf{Fungsi} \\
\midrule
\endhead

\bottomrule
\endlastfoot

Teks evidence & Isi potongan dokumen yang digunakan sebagai konteks downstream \\

Identitas dokumen & Mengidentifikasi dokumen pedoman asal evidence \\

Nomor halaman & Menunjukkan lokasi evidence pada dokumen sumber \\

Informasi metadata & Menyimpan atribut tambahan yang tersedia dari proses ingest \\

Identitas potongan & Membedakan potongan yang berasal dari dokumen atau halaman yang sama \\
\end{longtable}
\endgroup
